\documentclass[aspectratio=169]{beamer}
\usetheme{Madrid} % Un tema sobrio y muy académico
\usecolortheme{whale}

\title{Localización de una Fuente Sísmica}
\subtitle{Problema Inverso y Minimización Multivariable}
\author{Equipo de Trabajo}
\date{\today}

\begin{document}

\begin{frame}
    \titlepage
\end{frame}

\begin{frame}{Contenido}
    \tableofcontents
\end{frame}

\section{Problema Directo y Modelo de Atenuación}
\begin{frame}{El Modelo Geométrico y Atenuación}
    % Aquí colocaremos la definición de Ri y la ecuación de amplitud Azi.
\end{frame}

\section{Simulación y Ruido Gaussiano}
\begin{frame}{Generación de Datos Sintéticos}
    % Aquí explicaremos las coordenadas de los sensores y la inclusión del ruido.
\end{frame}

\section{Problema Inverso: La Función de Error}
\begin{frame}{Minimización por Mínimos Cuadrados}
    % La ecuación magistral E(x,y,z) y los resultados de Nelder-Mead.
\end{frame}

\section{Visualización y Curvas de Nivel}
\begin{frame}{Análisis Topológico del Error (z = -1.14)}
    % Espacio para integrar las gráficas de calor y concluir.
\end{frame}

\end{document}